\chapter{Swollen gums}
\index{Swollen gums}

\medskip
\noindent\textit{Educational reference; always seek professional advice for individual care.}

\section*{Definition and Overview}
Swollen gums, or gingival swelling, is the enlargement, puffiness, redness, or tenderness of the gums---the soft tissue that surrounds and supports the teeth. The gums are normally firm, pale pink (in people with lighter skin), and tightly bound to the teeth, and a change in their appearance or feel is a common reason for people to notice their oral health. The most frequent cause of swollen gums is inflammation from the accumulation of dental plaque, the sticky film of bacteria that forms on teeth every day; when plaque is not removed by brushing and flossing, it irritates the gums, producing the condition called gingivitis. Gingivitis is extremely common, reversible, and almost always preventable with good oral hygiene. In some cases, swelling can signal more significant problems, including periodontitis (inflammation spreading to the structures supporting the teeth), an abscess, a herpes-related mouth infection, reactions to medicines, hormonal changes in pregnancy, or rarely, serious disease. This chapter explains why gums swell, the common causes, how to tell the difference between minor and serious gum problems, how swollen gums are treated, and what can be done at home to prevent them. The central message is simple: gum health is a foundation of overall health, and regular dental care and good daily habits prevent the great majority of gum problems.

\section*{Epidemiology}
Gum problems are among the most common conditions in the world. Gingivitis, the mildest and commonest form of gum inflammation, affects a large proportion of adults at any given time and is near-universal at some point in life, and it is particularly common in people who smoke, have diabetes, or practise irregular oral hygiene. Periodontitis, the more serious form that damages the gums and the bone supporting the teeth, affects a significant minority of adults, becomes more common with age, and is a leading cause of tooth loss after the age of forty. Both conditions are more common in men and are strongly linked to social disadvantage, since access to dental care and resources for oral hygiene vary widely. Certain groups have higher rates of specific causes: pregnant women very commonly develop pregnancy gingivitis, peaking in the second trimester; people taking medicines such as calcium channel blockers for blood pressure, or phenytoin for epilepsy, can develop drug-related gum overgrowth; and adolescents are prone to pericoronitis, the inflammation of the gum flap around a partially erupted wisdom tooth. Tooth decay and dental abscesses are also extremely common and frequently present with a swollen gum around the affected tooth. In many countries, oral disease remains one of the most prevalent health problems, which highlights how much gum and dental health can be improved with prevention.

\section*{Aetiology and Pathophysiology}
The gums respond to irritation with the same inflammation seen elsewhere: blood vessels dilate, fluid and immune cells accumulate, and the tissue becomes red, swollen, and tender. The most important trigger is dental plaque, a biofilm of bacteria, food debris, and saliva that builds up on the teeth, especially along the gum line. If plaque is not removed within about two weeks, it hardens into calculus (tartar), which cannot be brushed away and keeps the gums chronically irritated. The severity depends on the bacteria present, the person's immune response, and modifying factors such as smoking, hormones, and medicines.

\begin{itemize}
    \item \textbf{Dental plaque and calculus:} the underlying cause of most gingivitis and periodontitis; plaque bacteria produce toxins that inflame the gums
    \item \textbf{Periodontitis:} inflammation spreading below the gum line into the pockets, destroying the ligament and bone that hold the tooth
    \item \textbf{Dental abscess and severe tooth decay:} infection tracks to the root, causing a swollen, painful gum and sometimes a visible boil
    \item \textbf{Pericoronitis:} inflammation of the gum flap over a partially erupted wisdom tooth, often with food trapping
    \item \textbf{Viral infections:} herpes gingivostomatitis causes painful, swollen gums with mouth ulcers, mainly in young children
    \item \textbf{Hormonal changes:} pregnancy and puberty increase the gum's response to plaque, producing swelling and bleeding that settles afterwards
    \item \textbf{Medications:} calcium channel blockers, phenytoin, and some immunosuppressants can cause drug-induced gum overgrowth
    \item \textbf{Nutritional deficiency:} vitamin C deficiency (scurvy) causes swollen, bleeding gums, though it is uncommon in many countries
    \item \textbf{Mouth breathing and dry mouth:} a dry mouth allows plaque to accumulate faster, and chronic mouth breathing inflames the front gums
    \item \textbf{Systemic disease:} some conditions, including diabetes, leukaemia, and certain autoimmune diseases, can affect the gums
\end{itemize}

The common thread is that plaque-related inflammation accounts for the great majority of cases, which is why the same prevention---thorough, regular cleaning of the teeth and gums---applies to almost everyone with swollen gums.
\section*{Clinical Features}
The features of swollen gums depend on the cause and how long it has been present.

\begin{itemize}
    \item \textbf{Redness and puffiness:} the gums look redder and fuller than usual, losing their firm, stippled appearance
    \item \textbf{Bleeding when brushing or flossing:} a hallmark of gingivitis; healthy gums do not bleed with normal cleaning
    \item \textbf{Tenderness:} gums may hurt when brushed or when eating hard or spicy food
    \item \textbf{Swelling around one tooth:} suggests an abscess or a local problem such as deep decay
    \item \textbf{A "boil" on the gum:} a small pus-filled swelling near the root of a tooth, which often comes and goes
    \item \textbf{Swelling of the gum flap at the back:} typical of pericoronitis, often with pain and difficulty opening the jaw fully
    \item \textbf{Bad breath and a bad taste:} from the bacterial build-up that accompanies gum inflammation
    \item \textbf{Receding gums or loose teeth:} later features of periodontitis, where the supporting structures have been destroyed
    \item \textbf{Painful ulcers with fever:} suggests viral gingivostomatitis, most often in children
    \item \textbf{Generalised overgrowth:} firm, enlarged gums covering more of the teeth, typical of drug-induced swelling
\end{itemize}

\section*{Differential Diagnosis}
Because gum swelling has many causes with different treatments, a dental or medical professional needs to identify the precise problem.

\begin{itemize}
    \item \textbf{Gingivitis:} inflammation limited to the gum margin, reversible with cleaning, and the commonest diagnosis
    \item \textbf{Periodontitis:} the same inflammation plus loss of the bone supporting the teeth, recognised on examination with a probe and on X-ray
    \item \textbf{Periapical abscess:} an infection at the root tip of a tooth, causing swelling near the root
    \item \textbf{Pericoronitis:} localised swelling of the gum flap over a wisdom tooth, distinct from generalised gum inflammation
    \item \textbf{Herpetic gingivostomatitis:} a painful primary herpes infection with ulcers, fever, and swollen gums, mostly in children
    \item \textbf{Drug-induced gum overgrowth:} firm enlargement in a person taking a culprit medicine, improving when the medicine is changed
    \item \textbf{Oral cancer and other growths:} any persistent, painless, or growing swelling or ulcer that does not heal within a few weeks warrants biopsy
    \item \textbf{Systemic causes:} diabetes, blood disorders, and autoimmune conditions can present with gum changes and need medical assessment
\end{itemize}

\section*{When to Seek Medical Care}
Many cases of swollen gums can be managed at home with improved brushing and flossing, and most people should start there. A dentist should be consulted for bleeding that does not settle with good hygiene, swelling that persists beyond a week or two, a painful swelling around a single tooth, a gum "boil", or suspected periodontitis with loose teeth or receding gums. A doctor should be involved if gum problems accompany unexplained fever, weight loss, or fatigue, and for children with painful mouth ulcers and fever. Anyone with diabetes should see their dentist regularly, since gum disease is more common in diabetes and can make blood sugar control harder. If gums swell without an obvious cause and persist despite good care, an assessment is worthwhile.

\textbf{Contact your local emergency services for:}
\begin{itemize}
    \item Swelling of the face, jaw, or throat, especially spreading or making breathing or swallowing difficult
    \item High fever with facial swelling, which may indicate a spreading dental infection
    \item Collapse or signs of a severe allergic reaction
\end{itemize}

\section*{Diagnosis and Investigations}
Gum problems are diagnosed mainly by examination, and imaging is added when the underlying bone or teeth need to be assessed.

\begin{itemize}
    \item \textbf{History:} the dentist asks about oral hygiene, bleeding, pain, medicines, pregnancy, smoking, diabetes, and duration
    \item \textbf{Clinical examination:} the gums are inspected and a periodontal probe measures the depth of the pockets around each tooth
    \item \textbf{Dental X-rays:} show the roots and supporting bone, distinguishing gingivitis from periodontitis and revealing abscesses and decay
    \item \textbf{Testing the teeth:} gentle tapping and cold testing identify a tooth with deep infection that may be the source of the swelling
    \item \textbf{Blood tests:} occasionally used to look for diabetes, vitamin deficiencies, or blood disorders when the pattern is unusual
    \item \textbf{Biopsy:} a persistent or suspicious swelling or ulcer is sampled to rule out oral cancer
\end{itemize}

\section*{Management}
Treatment depends on the cause, but the foundation for almost all gum problems is a thorough clean and improved daily hygiene.

\begin{itemize}
    \item \textbf{Professional cleaning (scale and polish):} removal of plaque and calculus by a dentist or hygienist, the definitive treatment for gingivitis
    \item \textbf{Improved home care:} brushing twice daily with a fluoride toothpaste, cleaning between the teeth daily, and using an antiseptic mouthwash if recommended
    \item \textbf{Periodontal therapy:} for periodontitis, deep cleaning of the root surfaces is performed, sometimes under local anaesthesia, and may be combined with surgery in advanced cases
    \item \textbf{Antibiotics:} used for a dental abscess, pericoronitis, or spreading infection, often alongside drainage or root canal treatment
    \item \textbf{Drainage of an abscess:} releasing the pus, either through the tooth, the gum, or the root, is the urgent step in relieving pain and controlling infection
    \item \textbf{Root canal treatment or extraction:} a tooth with deep infection is treated or removed, eliminating the source of the swelling
    \item \textbf{Medication review:} if drug-induced overgrowth is suspected, the prescribing doctor may adjust the medicine
    \item \textbf{Pregnancy gingivitis:} managed with meticulous hygiene and professional cleaning, settling after the birth
    \item \textbf{Symptom relief:} warm salt-water rinses, simple pain relief, and soft foods ease discomfort
\end{itemize}

\section*{Complications}
\begin{itemize}
    \item \textbf{Periodontitis and tooth loss:} untreated gum inflammation destroys the supporting bone, a leading cause of adult tooth loss
    \item \textbf{Spread of infection:} a dental abscess can spread into the jaw, face, or neck, and rarely the bloodstream
    \item \textbf{Chronic bad breath:} persistent bacterial build-up causes halitosis that is hard to mask
    \item \textbf{Effects on general health:} periodontitis is associated with a higher risk of cardiovascular disease and with poorer blood sugar control in diabetes
    \item \textbf{Impact on eating and nutrition:} painful, swollen gums can make chewing difficult
    \item \textbf{Psychological effects:} visible swelling, bleeding, or tooth loss can affect confidence
\end{itemize}

\section*{Prognosis and Outcomes}
The outlook for swollen gums is generally excellent, because the underlying problem is usually reversible gingivitis. With professional cleaning and improved daily care, gingivitis typically settles within a couple of weeks, and the gums return to a firm, pale, non-bleeding state. Periodontitis cannot be completely reversed, because bone that has been lost does not regrow, but it can be very effectively controlled: with regular professional care and good home hygiene, the disease is slowed or halted, teeth are retained for many years, and the gums stabilise. Abscesses and acute infections resolve quickly once drained and treated, and drug-induced and hormonal swelling improve with adjustments. The key predictor of a good outcome is early and sustained care: the sooner plaque is removed, the less damage is done. Regular dental check-ups, usually every six to twelve months, catch problems while they are still simple, and most people who follow that advice keep their natural teeth for life.

\section*{Prevention and Self-Care}
\begin{itemize}
    \item Brush teeth twice a day for two minutes with a fluoride toothpaste, reaching the gum line with gentle circular motions
    \item Clean between the teeth daily with floss or interdental brushes, since brushing alone misses about a third of the tooth surface
    \item See a dentist for a check-up and professional clean at least once a year, or more often if you have gum disease, diabetes, or smoke
    \item Stop smoking, which is one of the strongest risk factors for gum disease
    \item Manage diabetes well, since good blood sugar control reduces gum inflammation
    \item Use a soft-bristled brush and replace it every three months or when the bristles fray
    \item Rinse with warm salt water to ease temporary discomfort, but not in place of dental care
    \item Eat a balanced diet with plenty of fruit and vegetables, and limit sugary snacks and drinks
\end{itemize}

\section*{Sources and Further Reading}
The following organisations provide reliable information about gum health and dental care.

\begin{itemize}
    \item Australian Dental Association. \textit{Gum disease and oral health information.} https://www.ada.org.au/
    \item healthdirect Australia. \textit{Swollen or bleeding gums.} https://www.healthdirect.gov.au/
    \item NHS. \textit{Gum disease (gingivitis and periodontitis).} https://www.nhs.uk/
    \item Mayo Clinic. \textit{Swollen gums: causes and when to see a dentist.} https://www.mayoclinic.org/
    \item Australian Department of Health and Aged Care. \textit{Oral health and dental services in Australia.} https://www.health.gov.au/
    \item World Health Organization (WHO). \textit{Oral health information and guidelines.} https://www.who.int/
\end{itemize}

\clearpage
